% Laget automatisk av elevark/lag_elevark.py fra opplegg/magnetfelt-mobil.html. Endringer her blir overskrevet.
\documentclass[11pt]{article}
\usepackage{elevark}
\begin{document}
\elevarkhode{Fysikk 2 · Magnetiske felt og krefter}{Magnetfelt rundt en leder med mobilen}{https://lektorbo.no/opplegg/magnetfelt-mobil.html}{Opplegget på nett}

De fleste mobiler har et magnetometer som kan måle magnetfelt i tre retninger. Med appen phyphox kan dere måle feltet rundt en rett leder med strøm, og undersøke om det stemmer at $B = μ_{0}I/(2πr)$.

\section*{Hensikt}

\begin{itemize}
\item Måle magnetfeltet rundt en rett, strømførende leder i ulike avstander.
\item Undersøke om $B$ er omvendt proporsjonal med avstanden $r$.
\item Bestemme $μ_{0}$ fra målingene.
\end{itemize}

\section*{Utstyr}

\begin{itemize}
\item Mobil med appen phyphox (gratis) og verktøyet «Magnetometer».
\item Spenningskilde som kan gi 3–5~A likestrøm, helst med strømbegrensning.
\item Amperemeter, og en ledning som er minst 1~m lang (gjerne 1,5~mm² eller tykkere).
\item Linjal, tape og en liten magnet.
\end{itemize}

\section*{Framgangsmåte}

\begin{enumerate}
\item Finn sensoren: Åpne «Magnetometer» i phyphox og før en liten magnet over mobilen. Der utslaget er størst, sitter sensoren. Merk stedet med en tapebit på baksiden eller kanten.
\item Tape ledningen rett på bordet, og la de to tilførselsledningene gå langt bort fra målestedet. Legg mobilen flatt på bordet ved siden av ledningen, med sensormerket nærmest.
\item Når mobilen ligger i samme plan som ledningen, peker feltet fra ledningen rett opp eller ned ved sensoren. Les derfor av z-komponenten, $B_{z}$.
\item Still inn strømmen, for eksempel 5,0~A, og les av $B_{z}$. Snu så strømmen ved å bytte ledningene på spenningskilden, og les av $B_{z}$ igjen.
\item Feltet fra ledningen er halve differansen: $B = |B_{z,1} - B_{z,2}| / 2$. Da forsvinner jordas magnetfelt og andre felt som ikke endrer seg når strømmen snus.
\item Gjenta for avstander fra 1~cm til 6~cm mellom ledningen og sensoren. Hold mobilen i ro på samme sted på bordet, og flytt heller ledningen.
\end{enumerate}

\section*{Sikkerhet}

Ledningen blir varm med flere ampere. Slå av strømmen mellom målingene, og ikke kortslutt et batteri. Bruk en spenningskilde med strømbegrensning.

\section*{Oppgaver}

\begin{oppgaveliste}
\oppgave \fase{gjett}{Gjett}Hvor sterkt er feltet 2,0~cm fra en leder med 5,0~A? Sammenlign med jordas magnetfelt, som er omtrent 50~µT i Norge. Hvorfor er det lurt å snu strømmen i stedet for bare å slå den av og på?
\svar
\oppgave \fase{test}{Mål}Gjør målingene og fyll inn tabellen. Ligger punktene på en rett linje gjennom origo når dere plotter $B$ mot $1/r$? Hva betyr det?
\svar
\oppgave \fase{test}{Regn}Bruk stigningstallet til å bestemme $μ_{0}$. Hvor stort er avviket fra $4π \cdot 10^{-7}\,\mathrm{T\,m/A}$?
\svar
\oppgave \fase{forklar}{Forklar}Sensoren ligger noen millimeter inne i mobilen, altså litt over ledningen. Hvordan påvirker det målingene nær ledningen? Prøv å fjerne de nærmeste målingene og se hva som skjer med $μ_{0}$.
\svar
\oppgave \fase{forklar}{Vurder}Hvilke andre feilkilder er det i forsøket? Hvordan kunne dere gjort forsøket mer nøyaktig?
\svar
\oppgave \fase{utfordring}{Utfordring}Hold avstanden fast og varier strømmen fra 1~A til 5~A. Plott $B$ mot $I$. Hva forventer dere å se, og hva blir stigningstallet?
\svar
\end{oppgaveliste}

\end{document}
